 
 \chapter{Plus grand commun diviseur - plus petit commun multiple} 
 
 \begin{itemize}
 \item Calculer le plus grand commun diviseur (PGCD) des nombres suivants : 
 \begin{itemize}
 \it $168$ et $360$. 
 \it $252$ et $684$. 
 \it $336$ et $462$. 
 \it $1~840$ et $1~260$. 
 \it $18~150$ et $23~850$. 
 \it $33~390$ et $58~800$. 
 \it $315$, $819$ et $924$. 
 \it $252$, $693$ et $945$. 
 \it $2~520$, $3~150$ et $4~410$. 
 \it $7~560$, $10~080$ et $12~096$. 
 \end{itemize}
 \item Établir la liste des diviseurs communs aux nombres : 
 \begin{itemize}
 \it $4~200$ et $5~880$. 
 \it $1~440$ et $1~764$. 
 \it $3~780$, $4~320$ et $5~184$. 
 \it $10~584$, $11~520$ et $13~104$. 
 \end{itemize}
 \item Calculer le plus petit commun multiple (PPCM) et
 les trois multiples communs les plus simples de : 
 \begin{itemize}
 \it $360$ et $504$. 
 \it $252$ et $672$. 
 \it $972$ et $1~134$. 
 \it $720$ et $900$. 
 \it $168$, $252$ et $336$. 
 \it $120$, $180$ et $270$. 
 \end{itemize}
 \it \begin{enumerate}
 \item Calculer le PGCD et le PPCM des nombres $576$ et $1~080$. 
 \item Comparer le produit des deux résultats au produit des deux nombres. Énoncer le résultat obtenu.
 \end{enumerate}
 \it \begin{enumerate}
 \item Calcule le PGCD et le PPCM des nombres $99$ et $140$. 
 \item Quel est le PPCM de deux nombres premiers entre eux ? 
 \end{enumerate}
 \it \begin{enumerate}
 \item Calculer le PGCD et le PPCM des nombres $144$ et 
 $180$. 
 \item Que deviennent les résultats précédents lorsqu'on
 multiplie (ou divise) les deux nombres par $6$ ?
 Généraliser. 
 \end{enumerate}
 \it Démontrer que, si un nombre en divise deux autres, il divise leur somme, leur différence, et le reste de 
 leur division. 
 \it La division de deux nombres se fait exactement. Quel
 est leur PGCD et quel est leur PPCM ? Exemple : $6~375$
 et $375$. 
 \it Montrer que l'ensemble des diviseurs communs à deux nombres est le même que celui du plus petit de ces nombres et du reste de leur division. Que peut-on dire des PGCD ? \\ 
 \emph{Application.} — Remplacer la recherche du PGCD de $792$ et $240$ par la recherche du PGCD de deux nombres plus simples. Répéter cette opération afin d'obtenir un
 PGCD évident. (\emph{Méthode des divisions successives}.)
 \it Utiliser la méthode indiquée au numéro précédent pour la recherche du PGCD des nombres $2~021$ et $2~679$. 
 \it Trouver deux nombres non divisibles l'un par l'autre sachant que leur PGCD est égal à $336$ et leur somme 
 égale à $2~688$. 
 \it Par quel nombre inférieur à $100$ faut-il diviser $29~687$ et $35~312$ pour obtenir pour restes respectifs $47$ et $32$ ? 
 Quels sont les quotients ? 
 \it En divisant $809$ et $1~024$ chacun par un certain nombre, on trouve le même quotient et pour restes respectifs $27$ et $35$. Reconstituer les deux divisions. 
 \it On a planté des arbres également espacés sur le pourtour d'un terrain triangulaire dont les côtés mesurent $144$ m, $180$ m, et $240$ m. Sachant qu'il y a un arbre à chaque sommet et que la distance de deux arbres consécutifs est comprise entre $4$ mètres et $10$ mètres, calculer le nombres d'arbres plantés.
 \it Un ouvrier a touché pour trois mois successifs : $462$ F, $528$ F, et $594$ F. Trouver son salaire journalier sachant que c'est un nombre entier de francs 
 compris entre $20$ F et $30$ F. Trouver le nombre
 de jours de travail effectués chaque mois.
 \it On a fait carreler une pièce rectangulaire de $4,2$ m sur $2,24$ m. Sachant que les carreaux employés ont un côté compris entre $10$ cm et $25$ cm, calculer la longueur de leur côté et leur nombre. 
 \it On veut partager en coupons d'égale longueur quatre pièces d'étoffe mesurant respectivement $17,50$ m, $28$ m, $31,50$ m, et $42$ m. Trouver la plus grande longueur possible pour chaque coupon et le nombre total de ces coupons. 
 \it Deux règles égales de $504$ mm de longueur sont graduées l'une en $72$ parties et l'autre en $126$ parties. On fait coïncider leurs extrémités. Déterminer les traits de graduation qui coïncident. 
 \it Trouver les multiples communs à $6$, $8$ et $10$ compris entre $500$ et $1~000$. 
 \it Trouver les trois nombres les plus simples divisibles par les $10$ premiers nombres entiers. 
 \it Quel est le plus petit nombre qui donne $7$ pour reste quand on le divise par $12$, par $15$ ou par $16$ ? 
 \it Trouver un nombre qui donne $16$ pour reste quand on
 le divise par $24$ ou par $32$, et $8$ pour reste lorsqu'on le divise par $20$. 
 \it Trouver le plus petit nombre qui, divisé par $5$, $6$ ou $8$, donne respectivement pour restes $4$, $5$ et $7$. 
\it Deux cyclistes roulent dans le même sens sur une piste. Le premier fait un tour en $1$ minute $45$ seconde et le second en $1$ minute $36$ secondes. Sachant qu'ils sont partis ensemble de la ligne de départ, on demande après combien de temps passeront-ils ensemble cette ligne de départ. Combien de tours chacun d'eux aura-t-il effectués ? 
\it Des pavés rectangulaires qui ont $12$ cm de large et
$21$ cm de long ont servi à paver entièrement une place carrée. Calculer le côté de cette place sachant qu'il mesure un nombre entier de mètres compris entre $30$ mètres et $60$ mètres. 
\it Trois règles graduées de $960$ mm de long sont placées côté à côte de façon que leurs extrémités coïncident. Les divisions ont pour longueurs respectives $10$ mm, $12$ mm et $16$ mm. Déterminer les traits de divisions qui coïncident sur les trois règles. 
\it Une personne a acheté un certain nombre entier de mètres d'étoffe à $6,50$ F le mètre. Elle paye exactement avec des billets de $10$ F. Sachant que la dépense totale n'excède pas $200$ F, trouver le nombre de mètres achetés.
\it Un enfant compte ses timbres-poste par $12$, par $16$ et par $20$. Il lui en reste $8$ à chaque fois. En les comptant par $13$, il n'en reste plus. Combien possède-t-il de timbres ? 
\it En comptant les élèves d'une école par $9$, $10$ ou $12$, il en reste respectivement $8$, $9$ et $11$. En les comptant par $11$, il n'en reste pas. Trouver le nombre
d'élèves de l'école. 
\it Des autobus partent d'un même point dans quatre directions différentes. Les départs se font respectivement dans chaque direction toutes les $5$ minutes, $8$ minutes, $12$ minutes et $18$ minutes. Un départ simultané a lieu à $7$ heures le matin. Quelles sont les heures des autres départs simultanés de la journée ? 

 \end{itemize}